% Vietnamese translation for OI Public Library
% Source-File: IOI中国国家候选队论文/2003/侯启明/侯启明.doc
\documentclass[11pt]{article}
\usepackage{oipl-vn}

\title{Ứng dụng đơn giản của lý thuyết thông tin trong thi tin học}
\author{Hou Qiming}
\date{2003}

\begin{document}
\maketitle

\origtitle{Simple applications of information theory in informatics contests}
\sourcepath{IOI China candidate team papers / 2003 / Hou Qiming / Hou Qiming.doc}

\section{Mở đầu}

Lý thuyết thông tin là một nhánh rất quan trọng của khoa học máy tính. Tuy
trong các tài liệu về thi tin học trước đây hiếm khi thấy nội dung này, trên
thực tế lý thuyết thông tin có thể làm cơ sở để chứng minh tính tối ưu của
một số lời giải. Nếu vận dụng thích hợp trong thi đấu, nó thường đem lại hiệu
quả rất lớn. Vậy lý thuyết thông tin là gì?

\section{Sơ lược về lý thuyết thông tin}

Lý thuyết thông tin nghiên cứu bản chất của thông tin và quy luật truyền tải
thông tin. Trong thi lập trình, phần thường dùng nhất là quan hệ giữa độ bất
định và lượng thông tin. Nhờ nó, ta có thể thu được cận dưới tốt cho số bước
của một số bài toán tương tác; đây chính là điều một số tài liệu gọi là
\term{cận dưới theo lý thuyết thông tin}.

Ta nhắc lại vài khái niệm cốt lõi.

\paragraph{Định nghĩa.}
Nếu một biến ngẫu nhiên \(x\) có \(n\) giá trị có thể nhận, với xác suất lần
lượt là \(p_1,p_2,\ldots,p_n\), thì entropy của nó là
\[
  H(x)=f(p_1,p_2,\ldots,p_n)
  =-\sum_{i=1}^{n} C p_i\log p_i,
\]
trong đó \(C\) là một hằng số dương, thường lấy \(C=1\).

\paragraph{Định lý 1.}
Sau khi nhận được một thông tin có entropy \(h\) về biến ngẫu nhiên \(x\),
entropy của \(x\) giảm đi \(h\).

\paragraph{Định lý 2.}
Khi mọi giá trị có thể nhận của một biến ngẫu nhiên có xác suất bằng nhau,
entropy của biến đó là lớn nhất.

Các mệnh đề này trông rất gần với một số bài toán thi, đúng không? Ta xét vài
ví dụ cụ thể.

\section{Ví dụ 1: kiểm tra định lý 1}

Ngoài cửa sổ cầu thang giữa tầng hai và tầng ba ký túc xá của tác giả là mái
nhà của một căn phòng thấp bên cạnh. Ở đó có ba con mèo trắng: một con mèo
béo, một con mèo gầy có đuôi nguyên vẹn, và một con mèo gầy có đuôi không
nguyên vẹn. Khi trời lạnh, chúng thích nằm trên lò sưởi trong nhà. Vì vậy mỗi
con mèo có hai trạng thái: ở trong nhà hoặc ở ngoài trời.

Ba con mèo có tổng cộng \(8\) trạng thái, giả sử các trạng thái có xác suất
bằng nhau. Tác giả đứng ở cửa hàng nhỏ tầng một và chỉ có thể nhìn thấy số
lượng mèo đang ở trong nhà.

\paragraph{Câu hỏi 1.}
Xem toàn bộ trạng thái của ba con mèo là biến ngẫu nhiên \(x\). Khi chưa nhìn,
entropy của \(x\) là bao nhiêu?

Vì \(8\) trạng thái có xác suất bằng nhau,
\[
  H(x)=f\left(\frac18,\ldots,\frac18\right)=\log 8.
\]

\paragraph{Câu hỏi 2.}
Thông tin \(y\) thu được từ một lần nhìn có entropy bao nhiêu?

Số mèo trong nhà có thể là \(0,1,2,3\), với xác suất lần lượt
\[
  \frac18,\quad \frac38,\quad \frac38,\quad \frac18.
\]
Do đó
\[
  H(y)=
  -\left(
  \frac18\log\frac18+
  \frac38\log\frac38+
  \frac38\log\frac38+
  \frac18\log\frac18
  \right)
  =\log 8-6\log\frac38.
\]

\paragraph{Câu hỏi 3.}
Sau khi nhìn xong, entropy còn lại \(H'(x)\) là bao nhiêu?

Với các trường hợp số mèo trong nhà là \(0,1,2,3\), entropy điều kiện tương
ứng là
\[
  0,\quad \log 3,\quad \log 3,\quad 0.
\]
Vì vậy kỳ vọng entropy còn lại là
\[
  H'(x)=\frac18\cdot 0+\frac38\log 3+\frac38\log 3+\frac18\cdot 0
  =\frac68\log 3.
\]
Từ các công thức trên có thể kiểm tra được quan hệ
\[
  H(x)=H(y)+H'(x).
\]

\section{Ví dụ 2: Rods, IOI 2002}

Một \term{rod} là một thanh ngang hoặc dọc gồm ít nhất \(2\) ô vuông đơn vị
liền nhau. Trong một bảng \(N\times N\) có đặt hai rod, một ngang và một dọc.
Hai rod có thể có ô chung; khi có ô chung, ô đó được xem là đồng thời thuộc
cả hai rod.

Ban đầu ta không biết vị trí hai rod. Nhiệm vụ là dùng hàm thư viện
\(\operatorname{rect}(a,b,c,d)\) để xác định vị trí của chúng. Hàm này kiểm
tra hình chữ nhật \([a,b]\times[c,d]\), với \(a\le b\) và \(c\le d\). Nếu có
ít nhất một ô thuộc rod nằm trong hình chữ nhật đó, hàm trả về \(1\), ngược
lại trả về \(0\).

Khi thi, tác giả nhanh chóng nghĩ ra một phương pháp có số lần gọi nhiều nhất
là
\[
  6\log_2 n + C
\]
với \(C\) là hằng số. Vì con số này gần sát giới hạn bước, tác giả cố tối ưu
hệ số của \(\log_2 n\), nhưng không thu được gì; sau đó mới phát hiện điểm mất
điểm là xử lý thiếu trường hợp đặc biệt. Nhìn lại bằng lý thuyết thông tin,
ta thấy vì đề không cho xác suất, có thể giả sử mọi cách đặt rod là đồng xác
suất.

Gọi cách đặt rod là biến ngẫu nhiên \(x\). Nếu số cách đặt là \(f(n)\), thì
\[
  H(x)=\log f(n).
\]
Mỗi lần gọi \(\operatorname{rect}\) chỉ có hai kết quả, nên entropy lớn nhất
của câu trả lời là
\[
  H_{\max}(y)=\log 2.
\]
Do đó cận dưới theo lý thuyết thông tin cho số câu hỏi là
\[
  L=\frac{H(x)}{H_{\max}(y)}=\log_2 f(n).
\]

Theo nguồn gốc, trong bảng \(n\times n\), số cách đặt một rod là
\[
  n\binom{n+1}{2}=\frac{n^2(n+1)}2,
\]
và số cách đặt hai rod giao nhau là
\[
  \binom{n+2}{3}^2.
\]
Vì vậy
\[
  f(n)=
  \left(\frac{n^2(n+1)}2\right)^2
  -
  \left(\frac{(n+2)(n+1)n}{6}\right)^2
  =
  \frac{2n^6+3n^5-n^4-3n^3-n^2}{9}.
\]
Khi \(n\) đủ lớn,
\[
  L=\log_2 f(n)>\log_2\frac{2n^6}{9}
  \approx 6\log_2 n-2.2.
\]

Trong trường hợp xấu nhất, do các câu trả lời không nhất thiết luôn chia đôi
không gian trạng thái, số lần gọi thường không đạt đúng cận dưới, mà lệch
khoảng một hằng số. Vì vậy có thể xem \(6\log_2 n+C\) là cận dưới thực tế cho
số lần gọi lớn nhất của hàm \(\operatorname{rect}\). Khi đã có thuật toán đạt
dạng này, không cần mất thời gian vô ích để giảm hệ số; nên dành thời gian
kiểm tra trường hợp đặc biệt hoặc chuyển sang bài khác.

\section{Ví dụ 3: Coins}

Bài toán gốc từng được đề nghị cho một kỳ Olympic toán địa phương: có \(15\)
đồng xu, trong đó \(14\) đồng tốt và \(1\) đồng xấu. Mọi đồng tốt có cùng khối
lượng, còn đồng xấu có khối lượng khác. Biết trước một đồng là tốt, hãy dùng
cân hai đĩa để tìm đồng xấu trong không quá \(3\) lần cân.

Biến thể tổng quát là: có \(n\) đồng xu, trong đó \(n-1\) đồng tốt và \(1\)
đồng xấu; biết trước một đồng là tốt. Với \(k\) lần cân, có thể tìm ra đồng
xấu hay không? Nếu có thì in `POSSIBLE', ngược lại in `IMPOSSIBLE'.

Đánh số các đồng xu từ \(1\) đến \(n\), gọi \(x\) là số hiệu đồng xấu. Giả sử
mọi khả năng của \(x\) đồng xác suất, ta có
\[
  H(x)=\log n.
\]
Mỗi lần cân có ba kết quả: trái nặng hơn, phải nặng hơn, hoặc cân bằng. Vì vậy
\[
  H_{\max}(y)=\log 3.
\]
Do đó để tìm một đồng xấu trong \(n\) đồng cần ít nhất
\[
  \frac{H(x)}{H_{\max}(y)}=\log_3 n
\]
lần cân. Kết luận này chưa đủ để giải bài, nhưng nếu thiếu nó thì rất khó
chứng minh tính tối ưu của lời giải thực sự.

\subsection{Có đủ nhiều đồng tốt đã biết}

Giả sử có đủ nhiều đồng tốt đã biết. Gọi \(g(k)\) là số đồng nhiều nhất mà ta
có thể xử lý trong \(k\) lần cân để tìm một đồng xấu có khối lượng chưa biết.

Với \(k=1\), liệt kê trực tiếp cho thấy \(g(1)=2\).

Với \(k>1\), xét lần cân đầu tiên. Nếu kết quả cân bằng, gọi \(t\) là số đồng
chưa xác định không được đặt lên cân. Vì còn \(k-1\) lần cân để tìm đồng xấu
trong \(t\) đồng này, nên
\[
  t\le g(k-1).
\]
Nếu kết quả không cân bằng, đồng xấu nằm trong \(g(k)-t\) đồng đã lên cân. Nhờ
cận dưới theo lý thuyết thông tin, từ \(k-1\) lần cân còn lại ta suy ra
\[
  g(k)-t\le 3^{k-1}.
\]
Kết hợp hai bất đẳng thức:
\[
  g(k)\le g(k-1)+3^{k-1}.
\]

Cận trên này đạt được bằng cấu tạo. Lần đầu cân \(3^{k-1}\) đồng đang nghi
ngờ với \(3^{k-1}\) đồng tốt đã biết.

Nếu cân bằng, đồng xấu nằm trong phần còn lại gồm \(g(k-1)\) đồng, giải tiếp
theo định nghĩa của \(g\). Nếu bên các đồng tốt nhẹ hơn, đồng xấu nằm trong
\(3^{k-1}\) đồng nghi ngờ và là đồng nặng. Sau đó mỗi lần chia tập nghi ngờ
thành ba phần bằng nhau, cân hai phần; nếu cân bằng thì đồng xấu ở phần thứ
ba, nếu không thì ở phần nặng hơn. Trường hợp đồng xấu nhẹ tương tự.

Vì vậy
\[
  g(k)=g(k-1)+3^{k-1},
\]
và suy ra
\[
  g(k)=\frac{3^k+1}{2}.
\]

\subsection{Chỉ biết một đồng tốt}

Gọi \(f(k)\) là số đồng nhiều nhất mà ta có thể xử lý trong \(k\) lần cân khi
chỉ có một đồng tốt đã biết. Rõ ràng \(f(k)\le g(k)\). Nguồn gốc đưa ra một
cấu tạo để chứng minh thực ra \(f(k)=g(k)\).

Giả sử đồng \(1\) là đồng tốt đã biết. Lần cân đầu tiên cân các đồng
\[
  1,\ldots,\frac{3^{k-1}+1}{2}
\]
với các đồng
\[
  \frac{3^{k-1}+1}{2}+1,\ldots,3^{k-1}+1.
\]
Nếu cân bằng, dùng các đồng tốt vừa xác định để chuyển về tình huống có đủ
nhiều đồng tốt. Nếu không cân bằng, qua một lần cân phụ được thiết kế theo
các nhóm kích thước liên quan đến \(3^{k-2}\), hoặc chuyển về bài toán nhỏ
hơn với \(k-1\) lần cân, hoặc xác định rằng đồng xấu nằm trong một nhóm
\(3^{k-1}\) đồng và biết nó nặng hơn hay nhẹ hơn đồng tốt. Khi đó phần còn lại
giải bằng cách chia ba như trên.

Do đó với \(k\) lần cân, nhiều nhất có thể tìm đồng xấu trong
\[
  n=f(k)+1=\frac{3^k+3}{2}
\]
đồng xu.

\section{Kết luận}

Người đọc tinh ý sẽ thấy các ví dụ trong bài đều có thể giải mà không cần
kiến thức lý thuyết thông tin. Vậy lý thuyết thông tin có ý nghĩa gì ở đây?
Là một lý thuyết thuần túy, nó không nhất thiết là phương tiện trực tiếp để
giải từng bài cụ thể, mà là công cụ tổng quát để phân tích một lớp bài toán.
Nó có thể cung cấp cơ sở lý luận mạnh cho lời giải, và quan trọng hơn, giúp
ước lượng cận trên, cận dưới để định hướng cách xây dựng thuật toán.

Tóm lại, lý thuyết thông tin có rất nhiều đất dụng võ trong thi tin học.

\end{document}
